\scanpage{78}
i.e.
\[ \boxed{\sigma_0,\sigma_1\in\operatorname{Aff}(E)\qquad \sigma_0^2=\sigma_1^2=\operatorname{id}_E} \tag*{(2)} \]

(b)
\[ \boxed{r=(s_0,a_0)\in E\times\operatorname{Dr}(E),\quad(s_0,a_0)\ \text{incidents}} \tag*{(3)} \]

\uline{NB} Ces deux données définissent $\varphi_s,\varphi_a$ par
\[ \left\{\begin{aligned}\varphi_s(\underline{s}_n)&=s_n=u^n s_0\\
\varphi_a(\underline{a}_n)&=a_n=u^n a_0\end{aligned}\right. \tag*{\deleted{\unsure{(3)}}(4)} \]
avec
\[ u=\sigma_0\sigma_1 \tag*{\deleted{\unsure{(4)|(6)}}(5)} \]

\deleted{Mais on}
On trouve de $=$, pour $\underline r\in\underline{\operatorname{Rep}}(\Pi_\infty)$
\[ \underline r=\underline g\,\underline r_\infty\quad(\underline g\in\underline{\operatorname{Aut}}(\Pi_\infty)=D_\infty) \]
\[ \varphi_{\mathrm{rep}}(\underline r)=\deleted{\unsure{g}}\varphi_{\mathrm{rep}}(\underline g\,\underline r_\infty)=g\,\varphi_{\mathrm{rep}}(\underline r_\infty)=g.r. \tag*{\deleted{\unsure{(5)}}} \]
i.e.
\[ \varphi(\underline g\,\underline r_\infty)=g.r\qquad\text{si}\quad g=\varphi_G\,\underline g \tag*{(6)} \]
et les formules (4) \deleted{\unsure{(5)}} peuvent être
considérées comme déduites des \deleted{\unsure{fo}} (6)
par passage au quotient via
\diagram{\begin{tikzcd}[ampersand replacement=\&,column sep=small]
\& \operatorname{Rep}(\Pi_\infty) \arrow[dl,"\underline{\sigma}_1"'] \arrow[dr,"\underline{\sigma}_0"] \& \\
S(\Pi_\infty) \& \& A(\Pi_\infty)
\end{tikzcd}}
i.e. on doit avoir $\varphi_s,\varphi_a$ défini par
% archive sheet number: 77
